% Source: https://artofproblemsolving.com/community/c3h3164343_2020_emcc_individual_accuracy_test__exeter_math_club_competition
% Stable identity: aops-topic-3164343-q08
Write $7n=k^2+r$, where $k=\lfloor\sqrt{7n}\rfloor\le37$ and $r>0$. A zero remainder would force $7\mid n$, excluded. Squares modulo 7 are 0, 1, 2, 4, so $r=1,2$ are impossible. If $r\ge4$, the fractional part $\dfrac{r}{\sqrt{7n}+k}$ exceeds $\dfrac{4}{76}=\dfrac{1}{19}$. For $r=3$, one needs $k\equiv2$ or 5 modulo 7. The largest such $k\le37$ is 37, but it gives $n=196$, excluded as divisible by 7. The next is 33, giving $n=156$. Its fractional part is $\dfrac{3}{\sqrt{1092}+33}<\dfrac{1}{22}$, beating every $r\ge4$ case; among $r=3$ cases the fractional part decreases with $k$. Hence $\boxed{n=156}$.
